\documentclass[11pt]{ctexart}
\usepackage[a4paper,margin=2.3cm]{geometry}
\usepackage{booktabs}
\usepackage{amsmath}
\usepackage{amsthm}
\usepackage{amssymb}
\usepackage{stmaryrd}
\usepackage{tikz}
\usetikzlibrary{arrows.meta,positioning}
\usepackage{graphicx}
\usepackage[numbers,sort&compress]{natbib}
\usepackage{hyperref}
\hypersetup{colorlinks=true,linkcolor=blue!50!black,citecolor=blue!50!black}

\newtheorem{theorem}{定理}
\newtheorem{proposition}{命题}
\newtheorem{definition}{定义}
\newtheorem{lemma}{引理}

\setlength{\emergencystretch}{1em}
\linespread{1.15}

\title{\bfseries 智能体何时有资格宣称成功？\\[2pt]\large 面向开放世界结果的验证契约（Verification Contracts）}
\author{匿名作者\\[2pt]\small （本文为中文翻译稿，供内部审阅；术语首现处附英文原词）}
\date{}

\begin{document}
\maketitle

\begin{abstract}
\noindent 智能体正越来越多地执行这样一类任务：其完成状态无法从权威的程序状态中直接读取，而只能从截图、图像、文档、链接或一组上下文证据中\emph{推断}得到。现有系统常把这一判定交给某个模型裁判（model judge）。然而，即便一个判断既自信又在语义上正确，它也未必\emph{有资格}触发付款、访问授权、完成标记或工作流推进：验收规则可能在提交之后被改动，证据可能属于另一项任务，未消解的不确定性可能被强行压成一个二值答案，或者该次状态转移根本无法被事后重构。

我们研究的问题不同于"提升裁判准确率"：\emph{在何种条件下，一个概率性的观测才被允许授权一次确定性的状态转移？} 我们首先给出一个\emph{分离（separation）}结果——只要安全与不安全上下文在证据上不可区分，仅依赖证据（artifact-only）的裁判就不可能既在该等价类的安全成员上放行、又在该类上保持授权可靠（authorization-sound）——这就把问题置于"更好的分类器"力所不及的范围之外。随后我们提出\emph{验证契约（Verification Contract）}：一套将结果宣称、事前（ex-ante）证明义务、证据绑定、选择性验证、可容许转移与可重构回执联结起来的协议。我们定义四条可度量属性——事前一致性、证据绑定、选择性安全、决策可重构性。其中三条（事前一致性、证据绑定、可重构性）是对已承诺元数据的\emph{确定性}检查，因而是二值求值；只有一条\emph{概率性}谓词——选择性验证——才需要弃权值 $?$，以确保一个校准但不确定的观测不被强行压成二值。我们赋予这些谓词一套统一的三值授权语义与一个偏真守卫，并证明两条\emph{接口层}性质：一条\emph{违反传播不变量}——一个已知为 $\bot$ 的谓词值绝不会被已满足或未消解的同级谓词所掩盖（我们刻意不称其为“可靠性”，以免误读为经验安全宣称）——且守卫是\emph{弃权单调的}（消解一个未知量绝不会凭空制造出不安全放行）。我们明确指出：该不变量\emph{以"$\bot$ 被正确地产生"为前提}；它本身并不保证底层观测正确，因此它不是一项经验性的安全优势，而是一个组合式不变量，其现实价值由确证性研究来度量。本文核心的理论结果是与容量无关的\emph{上下文坍缩分离定理}；其余两条接口性质是描述守卫组合行为的辅助引理。

我们设计了 JOVE-Core：一个分组基准，由取得知情同意的数字基例（base case）与按属性对齐的匹配缺陷（matched defect）构成。其真值（ground truth）是一个\emph{独立评定}的下游危害——错误付款、错误授权、错误完成，或对合法提交者的错误拒绝——由对契约机制不知情的评定者判定，从而使指标不会退化为协议自身的规则。我们把确证端点分成两组分别预注册：\emph{结构一致性}端点（$P/B/R$，其真值由对冻结承诺的确定性重算给出，是机制的实现正确性检查）与\emph{外部结果}端点（不安全下游危害，由盲评者判定，承载经验优越性宣称）；只有后者可用于宣称相对基线的优越。外部结果的首要比较在预注册的 0.70 匹配自动化覆盖率上成对进行；完整风险—覆盖率曲线与 AURC 作为次要分析。其首要对照是"完整契约"对阵一个\emph{信息对齐（information-matched）}的裁判——后者获得完全相同的证据与观测载荷，但缺少结构性强制；带弃权接口的强策略感知裁判与确定性工具增强裁判作为次要参照。本文是一份遵循 Registered Report 模式的\emph{预注册研究手稿}：它预先冻结了问题、抽象、逐字段基线构造、目标样本量、统计功效、分析脚本与证伪标准。在确证数据被独立标注并分析之前，本文\emph{不作}任何经验性的优越性宣称。
\end{abstract}

\section{引言}
设想一个智能体，它要求用户在某截止时间前发布一段指定声明。用户提交了一张看似可信的截图，多模态裁判以 0.97 的置信度报告"成功"。系统应该放款吗？这个问题已不再是假想：新兴的智能体商务与智能体支付协议 \cite{coinbase2025x402,google2025ap2,anthropic2024mcp,ethereum2025erc8004} 允许自主智能体在结果被判定达成的那一刻即通过 HTTP 结算价值，于是裁判的"信念"正越来越直接地被接线到一次不可逆的转账上。

这个判断是不充分的。截图可能是旧的，可能属于另一场活动，也可能是在另一个策略版本下产生的。需求方可能在看到提交内容后才更改了标准。某个必需的服务方可能已经失效，而应用却仍在强求返回 PASS 或 FAIL。付款服务可能只保留了一个布尔值和一笔交易哈希，从而无从恢复"究竟是哪条规则、哪份证据授权了这次转账"。

这些并不都是分类错误。裁判完全可能\emph{正确地}识别出截图的内容，而该次转移\emph{依然}是不被允许的。真正缺失的对象，是"一次分类在什么条件下才被允许去控制状态"这一集合。

已有工作各自触及了这条边界的一部分。"LLM 充当裁判（LLM-as-a-Judge）"研究评测者的一致性与偏差 \cite{zheng2023judging,liu2023geval,wang2024fair}。选择性预测（selective prediction）研究模型何时应当弃权 \cite{geifman2017selective,geifman2019selectivenet,kamath2020selectiveqa,hendrickx2024reject}。媒体取证揭示伪造层面的风险信号 \cite{verdoliva2020media}。智能体基准评测由日益自主、拟人化的智能体所驱动的有后果的工作流 \cite{liu2023agentbench,zhou2023webarena,xie2024osworld,xu2024theagentcompany,yang2024sweagent,park2023generative}，而一条并行的安全研究线索则记录了\emph{为何}这些智能体自报的成功不能被照单全收：智能体会幻觉出结果 \cite{ji2023hallucination}、会依据被注入或被操纵的上下文行动 \cite{greshake2023injection}，并在面对逼真工具的沙盒中做出有后果却不安全的动作 \cite{ruan2024toolemu,kang2024exploiting}。人类计算（human computation）系统把人工劳动组织成可编程、有质量控制的资源，正对应了机器无法消解时所需的裁决环节 \cite{kittur2011crowdforge,little2010turkit,retelny2014flashteams,franklin2011crowddb}；而"验收标准随时间改变"本身即是一种概念漂移（concept drift）\cite{gama2014survey}。溯源与文档化实践记录血缘与适用范围 \cite{bender2018data,gebru2021datasheets,mitchell2019modelcards}。然而其中没有任何一项能\emph{单独}判定：不确定的外部证据在何时才有资格授权一次下游转移。

我们称之为\emph{开放世界结果授权（open-world outcome authorization）}问题。它源于一种不对称：观测始终是概率性的，而后果却是离散的。模型可以持有 0.73 的信念；付款系统却必须"放款或不放款"二选一。对某个分数设阈值，并不能解决策略漂移、跨任务重放、服务方失效或决策不可重构等问题。在区块链语境下，这一不对称被称为\emph{预言机问题（oracle problem）}——链下的、不确定的事实如何才能合法地授权一次链上的、确定的、不可逆的转移 \cite{caldarelli2020oracle,zhang2016towncrier,breidenbach2021chainlink}——而现有代币标准 \cite{vogelsteller2015erc20} 只标准化了转移本身，并未标准化其证据性前置条件。即便晚近面向智能体的链上标准已开始引入身份、声誉与验证注册表 \cite{ethereum2025erc8004}，它们登记的仍是\emph{参与方身份}与\emph{结果记录}，而非在给定冻结标准与任务绑定下、一个\emph{不确定}结果\emph{是否有资格}去授权一次转移——后者正是本文所填的空白。我们把开放世界结果授权视为该问题的一般形式，与任何具体账本无关。

我们提出\emph{验证契约（Verification Contract）}。它既不是法律合同，也不是对外部真相的形式化证明。它是一条\emph{可被机器强制执行}的授权边界：义务在证据到达之前就被承诺；证据被绑定到任务与版本；未消解状态可以让自动化停下；下游动作被列入白名单；且每一次决策都产生一份可重构的回执。

我们的贡献如下：
\begin{enumerate}
\item 我们形式化了开放世界结果授权问题，并证明一个\emph{分离}结果（定理~\ref{thm:separation}）：只要安全与不安全上下文共享完全相同的证据，任何仅依赖证据的裁判都不可能既在该歧义等价类上取得正放行覆盖、又在该类上保持授权可靠；因此该缺陷是\emph{信息层面与接口层面}的，而非统计层面的。
\item 我们定义了验证契约，以及一套带偏真守卫的三值授权语义——其中弃权值只被那条概率性谓词真正使用，其余三条谓词退化为确定性布尔检查。我们证明守卫的三条\emph{接口层}性质，并谨慎为它们命名以免夸大：一条\emph{违反传播不变量}（定理~\ref{thm:soundness}：一个已知为 $\bot$ 的值绝不会被掩盖——我们刻意不称其为“可靠性”）、\emph{弃权单调性}（命题~\ref{prop:abstain}：消解未知量绝不会制造不安全放行）、以及\emph{相对非冗余性}（定理~\ref{thm:necessity}：相对于所枚举的缺陷族，删去四条属性中的任意一条都会重新打开一个完整守卫本可阻止的匹配危害，并在命题~\ref{prop:completeness}中给出对所建模缺陷族的相对完备性）。这些是关于谓词取值的组合式不变量，而\emph{非}"谓词正确观测了外部真相"的保证；"一个 $\bot$ 是否在本应产生时被产生"这一经验问题留待确证研究。本文核心的\emph{理论}贡献是那条与容量无关的分离结果（定理~\ref{thm:separation}）；违反传播与相对非冗余性是刻画守卫接口行为的辅助引理，我们不把它们呈现为经验安全保证。
\item 我们设计 JOVE-Core，使每条属性都对应一个受控失败与一次消融（ablation），且其真值锚定于\emph{独立评定}的下游危害，而非契约自身的规则。
\item 我们预注册了一个信息对齐的首要基线、逐字段基线构造、风险—覆盖率评估、留出（held-out）泛化、人类重构实验、带功效分析的目标样本量，以及明确的证伪条件。
\end{enumerate}

我们用"开放世界"来命名这一\emph{观测范式}——结果从外部证据而非权威程序状态读取——而\emph{不是}宣称覆盖物理的、高风险的或对抗性不受约束的场景，后者不在确证范围之内。本文是一份遵循 Registered Report 模式的预注册研究手稿 \cite{nosek2018preregistration,chambers2022registered}。在没有确证证据之前，它不能宣称验证契约优于策略感知裁判。

\section{问题形式化}
设一个需求方或智能体声明结果 $q$，一个证据产出方提交 $e$，一个验证器获得观测 $o$，一个下游应用处于状态 $s$。该应用考虑执行一个正向的、有后果的动作 $a^+$，例如放款、授权访问、标记完成，或跨越一条不可逆的工作流边界。

传统裁判实现为
\begin{equation}
J(q,e)\rightarrow\{\mathrm{success},\mathrm{failure}\},
\end{equation}
可能附带置信度或理由。这一接口无法表达：规则是何时被固定的、证据属于哪项任务、必需的服务是否可用、未知如何传播、哪个动作被允许，以及决策如何被重构。

我们把\emph{不安全的正向转移（unsafe positive transition）}定义为：在未满足所声明的授权条件的情况下执行了 $a^+$。这一事件与"假阳性分类"相关，但并不等同：一个语义正确的观测仍可能是从另一项任务重放而来；一个真实的结果也仍可能是在"提交之后才被发明"的策略下被判定的。

\subsection{威胁模型}
产出方可能省略证据、重放过期材料、替换上下文、篡改文件，或利用评测者的歧义。需求方可能在提交后更改标准。服务方可能失效、漂移、意见相左或输出畸形。应用可能把超时、错误或弃权映射为成功。审计者可能拿不到私有的原始材料。

我们\emph{不}假设元数据诚实、模型误差相互独立、溯源完美，或人工复核无偏。我们\emph{不}覆盖每一种身份攻击、传感器攻陷、合谋策略或自适应媒体攻击。

\section{为何裁判准确率并不充分}
我们区分两类反例。前两类直接绑定到裁判的\emph{充分性}：即便是一个"神谕裁判"，仅凭证据也缺乏安全授权所需的信息。后两类暴露的是裁判周围的\emph{部署接口}；我们把它们陈述为接口反例、而非分类失败，因为一位细心的工程师无需改动裁判即可分别缓解它们。

\subsection{对裁判充分性的反例}
\paragraph{事后标准漂移。} 证据到达之后，起作用的验收规则改变了——例如新增了一个浏览量阈值。神谕裁判可以正确地套用旧规则或新规则，但没有任何判断能恢复"证据产生时究竟哪条规则在生效"。我们把它与合法的澄清区分开：需求方细化一个确实含糊的标准是良性的，契约通过对义务加版本来处理它，而非通过奖励拒绝。缺陷在于针对已冻结证据\emph{悄然改变了决定性规则}，而不在于任何提交后的沟通本身。

\paragraph{跨任务重放。} 一张截图确实包含了所需声明，却属于另一个账户、任务或策略版本。语义正确并不确立任务归属；裁判仅看到证据，原则上无法区分一次被绑定的提交与一次被重放的提交。

\subsection{对部署接口的反例}
\paragraph{被强制的二值输出。} 某个必需的服务方不可用，或某个强制模态不被支持，而应用接口却仍要求 PASS 或 FAIL。两种输出都没有有效的观测依据。这是答案接口的性质——只要接口允许弃权它就消失——因此它促成的是一套\emph{语义}，而非一个更准的裁判。

\paragraph{不可重构的转移。} 系统只存下了 PASS 与一笔交易哈希，却没有存策略版本、证据承诺、决定性义务、阈值或退化状态。这是一个日志与存储的选择，与裁判准确率正交；我们保留它，因为它在已部署的结算路径中普遍存在，也因为可重构性正是使其余三条可被审计的前提。

这些反例共同促成了裁判准确率之外的四条协议属性：两条因为裁判在原则上不充分，两条因为其周围的接口惯常性地丢弃了"安全地据判断行动"所需的信息。

\subsection{准确率与授权之间的一条分离}
前两个反例不只是"困难样本"；它们见证了一条任何分类改进都无法弥合的形式化缺口。把裁判建模为仅依赖所提交证据 $e$（可附带置信度）的任意函数 $J:\mathcal{A}\to\{\mathrm{success},\mathrm{failure}\}$，并令真实结果是请求 $q$、证据、以及\emph{上下文} $c$（任务身份、策略版本、溯源）的函数 $y^\star(q,e,c)$——而 $c$ 无法从 $e$ 恢复。裁判 $J$ 诱导出一条\emph{授权规则} $\mathrm{auth}_J$：当且仅当 $J$ 报告 $\mathrm{success}$ 时允许 $a^+$。

\begin{definition}[授权可靠性]
\label{def:authsound}
一条授权规则 $\mathrm{auth}_D$（决策过程 $D$ 连同其放行行为）在某任务族上是\emph{授权可靠（authorization-sound）}的，当且仅当对该族中的每个输入 $(q,e,c)$，它仅在 $y^\star(q,e,c)=\text{safe-to-authorize}$ 时才允许 $a^+$；等价地，它从不执行不安全的正向转移。可靠性只约束"放行"：一条永远拒绝的规则平凡地可靠却无用，故我们将该性质与覆盖率（被允许 $a^+$ 的输入占比）配对考察。
\end{definition}

\begin{definition}[上下文坍缩]
两个输入 $(q,e,c)$ 与 $(q,e,c')$ 对 $J$ 构成\emph{不可区分对}，若它们呈现相同证据 $e$，仅在上下文上不同，且 $y^\star(q,e,c)=\text{safe-to-authorize}$ 而 $y^\star(q,e,c')=\text{unsafe}$。
\end{definition}

\begin{theorem}[等价类条件下的可靠性--覆盖率分离]
\label{thm:separation}
设证据等价类 $E$ 至少包含一个可安全授权的输入和一个不安全输入。仅依赖证据的裁判 $J$ 不可能同时满足：（i）在 $E$ 的任一成员上允许 $a^+$；（ii）在 $E$ 上保持授权可靠。等价地，在 $E$ 上保持授权可靠，必然要求它对 $E$ 中安全成员的放行覆盖率为零。$J$ 在 $E$ 之外的准确率或覆盖率不改变这一局部不可能性。
\end{theorem}
\begin{proof}
由于 $J$ 仅读取 $e$（及由 $e$ 导出的量），它对 $E$ 中所有成员具有相同的放行行为。若它放行任一成员，就也会放行其中的不安全成员，从而违反授权可靠性；若它保持授权可靠，就必须拒绝 $E$ 的所有成员，包括安全成员，因而在 $E$ 的安全部分上放行覆盖率为零。$E$ 之外的决策不影响这一等价类条件下的权衡。
\end{proof}

跨任务重放与事后标准漂移各自实例化了这样一个对：相同证据 $e$，分别处于被绑定与被重放的上下文，或处于冻结与被篡改的标准之下。证据绑定谓词与事前一致性谓词通过让 $c$ 对守卫\emph{可见}来打破该对——这正是仅依赖证据的接口所丢弃的信息。因此定理~\ref{thm:separation} 一句话说明了：为何本文的贡献无法由一个更好的\emph{仅依赖证据}的裁判取得——缺陷在于信息与接口层面，而非统计层面。

\begin{figure}[t]
\centering
\resizebox{\columnwidth}{!}{%
\begin{tikzpicture}[
  font=\footnotesize,
  node distance=5mm and 8mm,
  art/.style={draw,rounded corners,align=center,inner sep=3pt,fill=black!4,minimum height=6mm},
  ctx/.style={draw,rounded corners,align=center,inner sep=3pt,minimum height=6mm},
  jud/.style={draw,align=center,inner sep=3pt,minimum height=6mm,fill=black!8},
  >={Stealth[round]}]
\node[art] (e) {相同证据 $e$\\（截图 / 摘要）};
\node[ctx,above right=4mm and 12mm of e,fill=green!8] (csafe) {上下文 $c$：已绑定\\$y^\star=$ 安全};
\node[ctx,below right=4mm and 12mm of e,fill=red!8] (cunsafe) {上下文 $c'$：重放/漂移\\$y^\star=$ 不安全};
\node[jud,right=26mm of e] (j) {仅依赖证据的 $J$\\$J(q,e,c)=J(q,e,c')$};
\draw[->] (e) -- (j);
\draw[->,dashed] (csafe) -- (j) node[midway,above,sloped]{\scriptsize 对 $J$ 不可见};
\draw[->,dashed] (cunsafe) -- (j) node[midway,below,sloped]{\scriptsize 对 $J$ 不可见};
\node[align=center,right=6mm of j,text width=26mm] (out) {单一固定输出 $\Rightarrow$ 在 $c$ \emph{或} $c'$ 上出错};
\draw[->] (j) -- (out);
\end{tikzpicture}%
}
\caption{定理~\ref{thm:separation} 背后的不可区分对。两个输入呈现\emph{相同}证据 $e$ 却有相反真值，仅在仅依赖证据的裁判无法观测的上下文 $c$（任务身份、策略版本、溯源）上不同。任何这样的 $J$ 在二者上必返回同一输出，故它要么在 $c'$ 上授权一次不安全转移，要么在 $c$ 上错误拒绝——这一缺口任何分类准确率都无法弥合。绑定谓词与事前谓词通过让 $c$ 对守卫可见来打破该对。}
\label{fig:indist}
\end{figure}

\paragraph{定理的适用范围，及其\emph{未}裁定之事。} 该定理约束的是仅读取证据的裁判；它正是"为何要把 $c$ 提供给决策"的动因。它有意对"通过某条通道被\emph{给定} $c$"的裁判——例如在提示词中以文本形式收到任务标识、版本与摘要——不作任何断言。这样的裁判在定理适用范围之外，我们并不宣称定理迫使它失败。这一点很重要，因为我们的首要确证基线（J1-IM，见 \S\ref{sec:study}）恰是这样一个被提供了信息的裁判。因此分离结果并未预先裁定确证比较；它只排除了"靠证据上的分类准确率来弥合缺口"这条路。把 $c$ \emph{作为提示文本提供给一个概率性裁判}是否与把它\emph{作为守卫中的确定性强制}提供同样安全，是一个独立的经验问题，而这正是确证研究要去回答、而非预设的东西。

\section{验证契约}
一个验证契约是
\begin{equation}
K=\langle I,P,B,V,A,R\rangle,
\end{equation}
其中 $I$ 含声明的结果与不可变约束；$P$ 含在证据到达前即承诺的证明义务；$B$ 把证据绑定到任务、策略版本与上下文；$V$ 是选择性验证器；$A$ 定义可容许的转移；$R$ 构造决策回执。

我们并不宣称这个元组是唯一可能的分解，但我们确实证明它是\emph{非冗余}的：在一个显式缺陷模型内（\S4.5，定理~\ref{thm:necessity}），去掉任意单条属性都会容许一次完整守卫本可阻止的有害转移，故无一属性可省。因此"最小性"分为我们严格区分的两层——一个模型内部的必要性定理，以及更强的经验主张（\emph{真实}缺陷确会实例化这些见证），后者正是 \S\ref{sec:study} 中每一项消融所要检验的。

\subsection{三值授权}
每个授权谓词在三值域 $\mathbb{T}=\{\top,\bot,?\}$ 中取值，其中 $?$ 表示未消解（未知）状态。我们采用 Kleene 强三值逻辑 \cite{kleene1952metamathematics,fitting1994kleene}，其合取 $\land$ 的相关子句为 $\top\land\top=\top$、对任意 $x$ 有 $\bot\land x=\bot$、以及 $\top\land{?}={?}$；数据库与监控中对未知值的处理用的正是这套 $\bot$ 占优语义，故我们复用而非另造。决定性的设计抉择在于：授权是通过一个\emph{偏真（truth-biased）}守卫 $\llbracket\cdot\rrbracket_\top$ 读出的——它只把 $\top$ 映射为\emph{放行}，而把 $\bot$ 与 $?$ 都坍缩为\emph{不放行}。于是"未知"永不为真：它无论经由何种组合规则，都不可能被提升为放行。

\paragraph{究竟哪些谓词真正用到第三个值。} 我们明确指出：四条谓词中只有一条在运行中是真正三值的。事前一致性（$P$，比较生效标准版本与已承诺版本）、证据绑定（$B$，比较证据摘要与任务/版本标识符与其承诺值）、可重构性（$R$，检查所需回执字段是否齐备）都是对已承诺元数据的\emph{确定性}比较：它们不携带校准分数、置信区间或保形集合，取值为 $\top$ 或 $\bot$，$?$ 仅保留给"输入确实缺失"这一退化情形。因此弃权值真正承重的只有选择性验证（$V$）——那唯一由概率模块支撑的谓词。我们保留统一的三值域，是因为守卫要对四者一致地组合，也因为当某个确定性输入缺失时 $?$ 仍须传播；但我们并不宣称对确定性谓词提出了新颖的三值处理。

\paragraph{把选择性验证拆成三级：观测 $O$、校准 $C$、裁决 $V$。} 审计指出“选择性验证”一词把三件应当分离的事捏在一起：从证据得到的\emph{原始观测}、把观测映为可比分数的\emph{校准}、以及把校准分数离散为 $\{\top,\bot,?\}$ 的\emph{裁决}。我们据此把 $V$ 显式分解为一条三级管线 $O\to C\to V$。\emph{观测层 $O$}：裁判或取证模块从证据产出一个原始信号（如逐义务的模型判断或风险分）。\emph{校准层 $C$}：把原始信号映为一个带校准保证的分数 $p\in[0,1]$，我们要求它是校准的 \cite{guo2017calibration}，并在可获得免分布覆盖保证时，于预注册风险水平 $\alpha$ 处暴露一个保形预测集 \cite{vovk2005algorithmic,angelopoulos2023conformal}；此层的质量由预注册的校准诊断（ECE 及可靠性图）\emph{独立}度量，见第 7 节。\emph{裁决层 $V$}：谓词值由事前固定的\emph{双阈值}规则给出：若校准下界超过接受阈值则 $\phi_V=\top$，若校准上界低于拒绝阈值则 $\phi_V=\bot$，否则 $\phi_V={?}$——即只要置信区间跨越决策边界，或保形集合不是单点集。这一分解使我们能把“观测错了”“校准坏了”“裁决阈值不当”三类失败归因到不同的层，而非笼统归给一个“$V$”。这正是不确定性\emph{被保留而非被四舍五入}之处：一个 $0.73$ 且区间很宽的分数产出 $?$（弃权），而非被强行取整为 $\top$。阈值与 $\alpha$ 在数据之前即冻结（见第~\ref{sec:study} 节）。我们强调它的边界：该离散化保证\emph{低裕度}或\emph{校准不良}的判断会以 $?$ 浮现，而非被取整成 $\top$；但它\emph{无法}阻止一个自信却\emph{错误}的 $\top$——一个落在接受区深处却事实错误的判断——通过守卫。这一残余风险不由语义处理，正是确证性风险—覆盖率评估要对照独立评定危害去度量的量。

一个正向动作只有在如下条件下才具资格：
\begin{align}
\mathrm{Authorize}(K,e,o,h,a^+)={}&\mathrm{Stable}(P)\land\mathrm{Bound}(B,e)\notag\\
&\land\mathrm{Accepted}(V,K,e,o,h)\notag\\
&\land\mathrm{Reconstructable}(R),
\end{align}
且 $a^+$ 被允许当且仅当 $\llbracket\mathrm{Authorize}(\cdot)\rrbracket_\top=\text{permit}$，即当且仅当 $\mathrm{Authorize}(\cdot)=\top$（图~\ref{fig:guard}）。由上文 $\bot$ 吸收、$?$ 保留的子句，单个 $\bot$ 合取项即拒绝，单个 $?$ 合取项（在无任何 $\bot$ 时）即弃权；因此一次放行要求\emph{全部四条}谓词各自独立地为 $\top$。缺失、冲突、不受支持或退化的证据会把相应谓词置为 $?$，并经偏真语义产出弃权而非强制二值。

\begin{figure}[t]
\centering
\resizebox{\columnwidth}{!}{%
\begin{tikzpicture}[
  font=\footnotesize,
  box/.style={draw,rounded corners,align=center,inner sep=2pt,minimum height=6mm},
  pred/.style={box,fill=black!3},
  outbox/.style={box,minimum width=13mm},
  >={Stealth[round]}]
\node[box] (x) {输入 $x=(K,e,o,h)$};
\node[pred,below left=8mm and 12mm of x] (p1) {$\phi_P$};
\node[pred,right=4mm of p1] (p2) {$\phi_B$};
\node[pred,right=4mm of p2] (p3) {$\phi_V$};
\node[pred,right=4mm of p3] (p4) {$\phi_R$};
\node[box,below=8mm of x,fill=black!5] (conj) {Kleene 合取 $\textstyle\bigwedge_i\phi_i\in\{\top,\bot,?\}$};
\node[box,below=6mm of conj] (guard) {偏真守卫 $\llbracket\cdot\rrbracket_\top$};
\node[outbox,below left=7mm and 6mm of guard,fill=green!8] (permit) {放行\\$a^+$};
\node[outbox,below=7mm of guard,fill=red!8] (reject) {拒绝};
\node[outbox,below right=7mm and 6mm of guard,fill=yellow!12] (abst) {弃权};
\draw[->] (x) -- (p1); \draw[->] (x) -- (p2);
\draw[->] (x) -- (p3); \draw[->] (x) -- (p4);
\draw[->] (p1) -- (conj); \draw[->] (p2) -- (conj);
\draw[->] (p3) -- (conj); \draw[->] (p4) -- (conj);
\draw[->] (conj) -- (guard);
\draw[->] (guard) -- node[left,pos=.55]{全为 $\top$} (permit);
\draw[->] (guard) -- node[right,pos=.55]{存在 $\bot$} (reject);
\draw[->] (guard) -- node[right,pos=.55]{存在 $?$、无 $\bot$} (abst);
\end{tikzpicture}%
}
\caption{作为决策系统的授权守卫。四条属性谓词在 $\{\top,\bot,?\}$ 中独立求值；其 Kleene 合取由偏真守卫读出——仅当全体一致为 $\top$ 时放行 $a^+$，出现任何 $\bot$ 即拒绝（定理~\ref{thm:soundness}），存在 $?$ 而无 $\bot$ 时弃权。删去任一谓词即删去一条通往合取的入边，也就删去了它对应的危害拦截块（定理~\ref{thm:necessity}）。}
\label{fig:guard}
\end{figure}

\subsection{违反传播不变量}
我们把该语义意在提供的\emph{唯一}保证隔离出来，并证明它作为一条接口性质成立，与裁判质量无关。我们刻意\emph{不}把它命名为“可靠性”（soundness），以免它被误读为一项经验性安全宣称；它准确的名字是一条\emph{违反传播不变量}：一个已知的违反一旦被某条谓词记为 $\bot$，就必然传播到守卫的输出而阻止放行。固定授权谓词集 $\Phi=\{\phi_1,\dots,\phi_4\}$（四条属性），令每个 $\phi_i:\mathcal{X}\to\mathbb{T}$ 是输入 $x=(K,e,o,h)$ 上的部分求值。称正向转移 $a^+$ 在 $x$ 处被\emph{不安全授权}，若 $a^+$ 被放行时存在某条声明义务在 $x$ 处已知被违反，即 $\exists i:\phi_i(x)=\bot$。

\begin{theorem}[违反传播不变量]
\label{thm:soundness}
在偏真三值守卫下，对每个输入 $x$、以及每个把 $a^+$ 门控于 $\llbracket\bigwedge_i\phi_i(x)\rrbracket_\top=\text{permit}$ 的应用，不会发生任何被不安全授权的转移。等价地，一条已知的义务违反绝不会被未消解或已满足的同级谓词所掩盖。
\end{theorem}
\begin{proof}
设对某 $j$ 有 $\phi_j(x)=\bot$。由 $\bot$ 吸收子句，无论其余合取项在 $\{\top,?\}$ 中取何值，均有 $\bigwedge_i\phi_i(x)=\bot$。偏真守卫将 $\bot\mapsto\text{不放行}$，故 $a^+$ 不被放行，与"被不安全授权"的定义矛盾。故不存在这样的转移。
\end{proof}

\paragraph{该不变量所声称与\emph{未}声称之事。} 我们刻意陈述定理~\ref{thm:soundness} 的窄义，以免它被误读为经验性安全宣称——这也是我们避用“可靠性”一词的原因。它\emph{恰好}证明了：\emph{若}某条谓词取值为 $\bot$（一条已知违反），\emph{则}守卫不会放行；即义务被违反的知识绝不会被组合规则冲掉。它\emph{没有}证明谓词正确反映了外部真相。若一个取证模块把一次不安全情形误标为 $\top$（假阴性漏检），守卫将会放行——这不违反该不变量，因为前件（某个 $\bot$）从未出现。因此该定理是关于\emph{$\bot$ 一旦出现如何传播}的组合式保证，而非关于\emph{$\bot$ 是否会在应当出现时出现}的经验保证。前者靠证明确立；后者是第~\ref{sec:study} 节确证研究的核心。这一区分正是本文不预称经验优越性的原因。

\begin{proposition}[弃权单调性]
\label{prop:abstain}
将任意谓词从 $?$ 消解到 $\{\top,\bot\}$ 中的某个值，绝不会把一次拒绝或弃权翻转为不安全放行。具体地，若 $\bigwedge_i\phi_i(x)\in\{\bot\}$，则把某个 $\phi_k$ 由 $?$ 精化后仍保持结论为 $\bot$，除非该 $\bot$ 本就来自 $\phi_k$；且只有当\emph{所有}谓词求值为 $\top$ 时才可能放行。
\end{proposition}
\begin{proof}
放行要求 $\bigwedge_i\phi_i(x)=\top$，这当且仅当每个 $\phi_i(x)=\top$。若某 $\phi_j(x)=\bot$，则由 $\bot$ 吸收，无论把某个处于 $?$ 的 $\phi_k$ 精化为何值，合取仍为 $\bot$。若合取为 $?$（无 $\bot$、至少一个 $?$），则把某个 $?$ 消解为 $\bot$ 得到 $\bot$（拒绝），消解为 $\top$ 则要么仍余其他 $?$（仍弃权），要么使全体为 $\top$（放行）——而这只在不再有已知违反时发生。任何情形下都不会由未知的消解\emph{凭空制造}出一次绕过某个 $\bot$ 的放行。
\end{proof}

\subsection{契约属性}
\textbf{事前一致性（ex-ante consistency）}要求与决策相关的义务在第一份证据被承诺之前即被冻结。任何修订都会生成一个新的不可变版本。

\textbf{证据绑定（evidence binding）}要求被接受的证据被承诺到任务身份、策略版本、摘要与所需上下文。摘要确立身份，而非真值。这条义务是内容溯源标准（如 C2PA \cite{c2pa2023spec}）与可验证凭证 \cite{sporny2022vc} 在授权时刻的对应物——它们同样把证据绑定到签发者与上下文；我们的绑定谓词在有溯源时消费之，但不假设它一定存在或诚实。

\textbf{选择性安全（selective safety）}要求：当依赖不可用、存在硬冲突、强制输入不受支持，或不确定性未消解时，必须弃权。下游代码不得把弃权重新解释为接受。

\textbf{决策可重构性（decision reconstructability）}要求一位独立审计者能恢复起支配作用的策略、证据承诺、决定性观测、验证器状态与被允许的转移。

\paragraph{强制不变量（属于框定，非贡献）。} 若某应用把 $\mathrm{Authorize}=\top$ 强制为每个 $a^+$ 的必要前提，则一个显式为假的契约谓词无法被该合规应用悄然绕过。这是一条同义反复的接口性质——强制一个守卫即意味着该守卫被强制——陈述它只为固定术语。它不是关于外部真相、策略正当性或语义可靠性的定理，我们不把它列入贡献。

\subsection{四条属性的相对非冗余性}
\S4.1 未以断言方式声称最小性；我们也\emph{不}宣称一般意义上的“最小性”。我们现在在一个显式缺陷模型内证明一个更窄的性质——\emph{相对非冗余性}：相对于所枚举的缺陷族，没有任何属性是冗余的，因为去掉任意一条都会容许一次完整守卫本可拦截的有害转移。该定理是一个\emph{关于守卫在所建模缺陷族上}的陈述；真实缺陷是否实例化这些见证是我们的消融所检验的独立经验问题，二者严格分开。它\emph{不}宣称四条属性对一切可能缺陷是最小充分集，只宣称在所建模族内每一条各自不可省。

我们先在定义~\ref{def:harm} 中命名守卫存在所要防止的危害。令一个\emph{世界} $w$ 为输入 $x=(K,e,o,h)$ 固定下游后果所依赖的基本事实：证据时刻生效的标准是否接受 $e$、$e$ 是否绑定到本任务与版本、强制观测是否被有效消解、以及该转移事后是否可重构。

\begin{definition}[授权危害与审计危害]
\label{def:harm}
$(x,w)$ 处一个被放行的转移 $a^+$ 是一次\emph{授权危害} $H_{\mathrm{auth}}$，若 $w$ 使 $a^+$ 在一条并不成立的标准或任务绑定下转移价值或授予访问，或错误地拒绝了一位合法提交者。它是一次\emph{审计危害} $H_{\mathrm{audit}}$，若 $a^+$ 被放行，但没有任何审计者能仅凭回执恢复起支配作用的策略、证据承诺、决定性观测与被允许的转移。
\end{definition}

\begin{definition}[属性匹配的缺陷族]
\label{def:defects}
对每条属性 $i\in\{P,B,V,R\}$，令 $D_i$ 是恰好呈现其匹配缺陷、且无其他缺陷的输入集：$D_P$，生效标准不同于已承诺标准（事后漂移）；$D_B$，$e$ 格式良好但绑定到另一任务或版本（重放）；$D_V$，某强制观测不可用、不受支持或硬冲突（强制二值）；$D_R$，回执遗漏了一个重构决策所需的字段。每个 $\phi_i$ 是这样一个谓词：它在 $D_i$ 上取值为 $\bot$（对 $P,B$）、在 $D_V$ 上取值为 $?$、并对 $R$ 在可重构性未确立时取 $\bot$。
\end{definition}

记完整守卫为 $G=\llbracket\bigwedge_{i\in\{P,B,V,R\}}\phi_i\rrbracket_\top$，去掉属性 $i$ 的守卫为 $G_{-i}=\llbracket\bigwedge_{j\neq i}\phi_j\rrbracket_\top$。

\begin{theorem}[相对非冗余性]
\label{thm:necessity}
对每条属性 $i\in\{P,B,V,R\}$，都存在一个输入 $x_i$，在其上 $G$ 不放行 $a^+$ 而 $G_{-i}$ 放行，且在其上放行 $a^+$ 对 $i\in\{P,B\}$ 是一次 $H_{\mathrm{auth}}$ 危害、对 $i=V$ 是一次未消解（强制二值）危害、对 $i=R$ 是一次 $H_{\mathrm{audit}}$ 危害。故没有任何属性冗余：每一条都严格扩大被拦截的有害转移集合。
\end{theorem}
\begin{proof}
固定 $i$，取 $x_i\in D_i$ 且其余缺陷均不存在，于是对 $j\neq i$ 有 $\phi_j(x_i)=\top$，并由定义~\ref{def:defects} 有 $\phi_i(x_i)\in\{\bot,?\}$。则 $G_{-i}(x_i)=\llbracket\bigwedge_{j\neq i}\top\rrbracket_\top=\text{permit}$，即缩减守卫授权了 $a^+$。对完整守卫，$\bigwedge_j\phi_j(x_i)$ 在位置 $i$ 有一个 $\bot$ 或 $?$ 合取项；由 $\bot$ 吸收与 $?$ 保留子句，它等于 $\bot$ 或 $?$，偏真守卫将二者皆映射为不放行，故 $G(x_i)$ 不放行。$G_{-i}$ 下被放行的转移在所述意义上有害：在 $D_P$/$D_B$ 上世界使 $a^+$ 在不成立的标准/绑定下行动（$H_{\mathrm{auth}}$）；在 $D_V$ 上它在无有效观测依据时强推裁决；在 $D_R$ 上它提交了无审计者能重构的转移（$H_{\mathrm{audit}}$）。于是对每个 $i$，去掉 $i$ 都会容许一次完整守卫本可拦截的危害。
\end{proof}

定理~\ref{thm:necessity} 是我们对“最小性”直觉所愿意形式化的\emph{全部}内容，其正名为\emph{相对非冗余性}：四条谓词是逻辑上相互独立的守卫，各自是其匹配危害在模型内的唯一拦截者，故删去任一都会重开一次有害转移。它\emph{并未}宣称四者对 $\bigcup_i D_i$ 之外的危害联合充分，也不宣称不存在某个更小的替代分解；这一点在下文界定。

表~\ref{tab:defects} 沿着"把此必要性结果与分离定理相连"的轴线陈列该族：每个缺陷所依赖的\emph{判别性事实}。对 $P$ 与 $B$，该事实是一段已承诺的上下文——生效标准版本、证据被捕获时对应的任务/版本——它按构造\emph{不在证据中}，因而对仅依赖证据的裁判不可见（定理~\ref{thm:separation}）；守卫靠对冻结承诺的确定性比较来消解它，而非靠推断。对 $V$，该事实是某强制观测的\emph{可得性}本身，二值裁判无法表达它，只能把它捏造成一次强制裁决；弃权值正是使之可被表示的东西。对 $R$，该事实是事后的：记录是否足以支撑重构，这是任何裁决时刻的置信度都无法提供的。自上而下读该表：两行取 $\bot$ 的正是分离定理预言仅依赖证据的裁判无法闭合的情形，取 $?$ 的一行是二值裁判甚至无法命名的情形，而 $R$ 行则完全与裁判正交——因此这四条不是一份任意清单，而是对"一次授权可能出错、且概率裁决单独无法覆盖"之诸方式的一个划分。

\begin{table}[t]
\centering
\footnotesize
\caption{定义~\ref{def:defects} 的属性匹配缺陷族，按每个缺陷所依赖的事实组织。末列陈述为何一个概率性、仅依赖证据的裁判在原则上无法闭合此案——把必要性结果（定理~\ref{thm:necessity}）与分离结果（定理~\ref{thm:separation}）相连。}
\label{tab:defects}
\resizebox{\columnwidth}{!}{%
\begin{tabular}{@{}llllc@{}}
\toprule
属性 & 缺陷 $D_i$ & 判别性事实 & 守卫如何消解 & \shortstack{在 $D_i$ 上的\\谓词值} \\
\midrule
$P$（事前标准）& 事后标准漂移 & 生效 vs.\ 已承诺标准版本 & 与冻结承诺作相等比较 & $\bot$ \\
$B$（证据绑定）& 跨任务 / 版本重放 & 证据被绑定到的任务+版本 & 与已承诺绑定作相等比较 & $\bot$ \\
$V$（选择性裁决）& 缺失/冲突观测下的强制二值 & 强制观测是否可消解 & 产出 $?$；弃权而非猜测 & $?$ \\
$R$（可重构性）& 回执遗漏决定性字段 & 记录是否支撑事后审计 & 要求回执字段齐备 & $\bot$ \\
\bottomrule
\end{tabular}%
}
\end{table}

\begin{proposition}[对所建模缺陷的相对完备性]
\label{prop:completeness}
设确证范围内每个不安全或不可审计的输入都落在 $\bigcup_{i}D_i$ 中（每个负例至少带一个匹配缺陷）。则 $G$ 不放行任何 $H_{\mathrm{auth}}$ 或 $H_{\mathrm{audit}}$ 转移：守卫是可靠的（定理~\ref{thm:soundness}），且相对于所建模缺陷族是完备的。
\end{proposition}
\begin{proof}
令 $x$ 不安全或不可审计。由假设 $x\in D_i$ 对某 $i$ 成立，故 $\phi_i(x)\in\{\bot,?\}$；由吸收/保留有 $\bigwedge_j\phi_j(x)\neq\top$，守卫不放行。反之任何被放行的 $x$ 对所有 $j$ 有 $\phi_j(x)=\top$，故 $x\notin\bigcup_iD_i$，在范围内既不不安全也非不可审计。于是"被放行" $=$ "在所建模族内安全且可审计"。
\end{proof}

我们把命题~\ref{prop:completeness} 的假设陈述为一个\emph{范围条件}，而非关于世界的事实：一个自适应对手可以构造 $\bigcup_iD_i$ 之外的危害，命题在那里不作任何断言。这恰是威胁模型所声明、并在局限性中重访的边界，现在被精确化：完备性仅相对于所枚举的、由研究者实例化的缺陷族被断言，而确证研究的留出类与变异源正检验该相对保证能否在分布漂移下存续。

\section{与相邻抽象的关系}
软件契约在机器可观测状态上规定条件 \cite{meyer1997object}；这里后置条件本身必须从不完整证据中推断。携带证明的代码（proof-carrying code）附带一份机器可校验的安全证明 \cite{necula1997proof}；开放世界的截图与上下文宣称通常不容许这样的证明。运行时验证根据形式化属性检查轨迹 \cite{leucker2009brief}；它本身并不为一个外部结果定义证据充分性。访问控制与授权模型——从 Saltzer 与 Schroeder 的经典保护原则 \cite{saltzer1975protection}、基于能力的组合 \cite{miller2006capability}，到现代的策略即代码引擎 \cite{opa2023}——在受信任、机器可观测的主客体上裁定权限；开放世界情形的不同恰在于：决定性事实（结果是否发生、在何标准下）本身是不确定且被推断的，故策略必须门控于一个概率性观测、而非已认证状态。选择性预测加入了弃权，但没有事前承诺、任务绑定或下游强制。溯源记录血缘，而一份完整的血缘记录未必解释"是哪条义务授权了转移"。最后，人机依赖文献表明，原始模型置信度常令人类信任失准，且人机组合团队并不可靠地胜过单独的更强一方 \cite{schemmer2023reliance,vaccaro2024combinations}；我们的人类研究（第~\ref{sec:study} 节）正是围绕这一风险设计，度量的是不当依赖，而非假设结构化证据一定有帮助。

\paragraph{完整性的密码学保证 vs.\ 结果的授权。} 一位安全读者会问：为何这个守卫不就是三种成熟完整性机制的一个实例？区块链预言机——从 Town Crier 的认证数据源 \cite{zhang2016towncrier}、DECO 的 TLS-预言机证明 \cite{zhang2020deco}，到生产级预言机网络 \cite{ellis2017chainlink} 与兼容传统体系的去中心化身份 \cite{maram2021candid}——以及更广的预言机问题 \cite{caldarelli2020oracle}，认证的是一个链下数据的\emph{来源}：它们证明某个值确实源自给定服务器或 TLS 会话。可信执行环境 \cite{costan2016sgx,intel2021tdx,sabt2015tee} 证明某项指定\emph{计算}在飞地内未经修改地运行。零知识与可验证计算系统 \cite{goldwasser1989knowledge,parno2013pinocchio,ben2014succinct,thaler2022proofs} 证明一个被承诺的输入被一个固定函数正确处理。三者都是验证契约的补充而非替代，原因是结构性的：每一种都认证\emph{一条通道或一次计算的完整性}，而开放世界的问题是\emph{一个不确定结果的资格}。一张 TLS 认证的截图、一个飞地证明的裁判、一次 ZK 证明的分类，都与 \S3 的四类缺陷完全相容：证据可以是真实取得的、裁判可以未经篡改地运行、证明可以有效，而与此同时标准悄然漂移、证据从另一任务重放、某强制服务方失效、或回执不可重构。这些机制回答的是"这串字节是否被忠实地取得并处理？"；它们并不回答"这个结果，在这条冻结标准与这个绑定下，是否被允许去改变状态？"第四条相邻线索——结算终局性，即一次转移在某共识规则下变得不可逆的那一点，无论是概率性的 \cite{nakamoto2008bitcoin,garay2015backbone} 还是可证终局的 \cite{buterin2020gasper,buchman2016tendermint}——同样正交：它支配的是一次状态改变\emph{何时}不再能被撤销，而非证据\emph{是否}曾授权它；也正因结算不可逆，先于它的授权决策才必须在一条冻结标准下作出。守卫把这些机制的输出当作\emph{更强的证据通道}来消费——一个被认证的 $V$、一个溯源认证的 $B$——但它自身的贡献，即把放行锚定到外部下游危害而非通道完整性，与它们全都正交（表~\ref{tab:positioning}）。它们有据可查的失效模式——例如攻破飞地机密性的瞬态执行攻击 \cite{vanbulck2018foreshadow}——本身正是我们把这些当作可加固输入、而非授权边界的原因。

无任何组件被宣称为单独新颖。贡献在于那条联结事前承诺、绑定、选择性决策与可重构转移的授权边界，连同一套属性对齐的评估——其真值是外部下游危害，而非边界自身的规则。

表~\ref{tab:positioning} 沿授权边界必须覆盖的五个维度把差距具体化。每条相邻工作各解决其一子集——选择性预测贡献弃权，溯源贡献绑定，策略即代码与软件契约在\emph{受信任机器状态}上贡献事前承诺——但无一者同时做到：在证据之前即承诺验收规则、把一个\emph{不确定}证据绑定到任务与版本、在观测未消解时弃权、把评估锚定到外部危害而非其自身裁决、并产出可重构回执。末行是唯一被完全填满的一行；正是这一模式、而非任何单元格的新颖性，构成贡献。

\begin{table}[t]
\centering
\footnotesize
\caption{与相邻抽象的定位对比。$\bullet$：作为设计目标直接提供该属性；$\circ$：部分或仅附带地提供（例如仅在受信任、机器可观测的状态上，或作为血缘之类的副产品）；---：未处理。“真值 $=$ 外部危害”询问的是：该方法的正确性是否对照其自身决策规则\emph{之外}的一个后果来判定。}
\label{tab:positioning}
\resizebox{\columnwidth}{!}{%
\begin{tabular}{@{}lccccc@{}}
\toprule
方法 & \shortstack{事前规则\\承诺} & \shortstack{不确定证据的\\任务/版本绑定} & \shortstack{选择性\\弃权} & \shortstack{真值 $=$\\外部危害} & \shortstack{可重构\\回执} \\
\midrule
LLM 充当裁判 \cite{zheng2023judging,liu2023geval} & --- & --- & $\circ$ & --- & --- \\
选择性预测 \cite{geifman2017selective,hendrickx2024reject} & --- & --- & $\bullet$ & --- & --- \\
媒体取证 \cite{verdoliva2020media} & --- & $\circ$ & --- & --- & --- \\
内容溯源 / VC \cite{c2pa2023spec,sporny2022vc} & --- & $\bullet$ & --- & --- & $\circ$ \\
区块链预言机 / zkTLS \cite{zhang2016towncrier,zhang2020deco} & --- & $\circ$ & --- & --- & $\circ$ \\
可信执行环境 \cite{costan2016sgx,intel2021tdx} & --- & $\circ$ & --- & --- & $\circ$ \\
ZK / 可验证计算 \cite{parno2013pinocchio,ben2014succinct} & $\circ$ & $\circ$ & --- & --- & $\circ$ \\
访问控制 / 策略即代码 \cite{saltzer1975protection,opa2023} & $\circ$ & $\circ$ & --- & --- & $\circ$ \\
软件契约 / PCC / RV \cite{meyer1997object,necula1997proof,leucker2009brief} & $\circ$ & --- & --- & --- & $\circ$ \\
\midrule
\textbf{验证契约（本文）} & $\bullet$ & $\bullet$ & $\bullet$ & $\bullet$ & $\bullet$ \\
\bottomrule
\end{tabular}%
}
\end{table}

\paragraph{定位。} 本文不改进任何模型；它规定并评估一套包裹裁判的授权协议。因此它的自然归属是可信机器学习评估与系统安全的交汇处——那些关注策略强制、溯源与可审计结算的场所——而非一个建模赛道。我们明确指出这一点，以便确证宣称被作为一项协议与评估贡献、而非一项学习成果来评判。

\section{契约治理下的验证}
原型把一段自然语言请求编译为带类型的义务，规定证据模态、强制状态、决策方法与失败行为。一个模型辅助的编译器与确定性回退相配对，使得服务方失效不能移除策略边界。

被提交的证据被绑定到任务标识符、策略版本、服务端回执时间与证据摘要。观测模块暴露确定性规则检查、有限的取证风险信号，以及逐义务的多模态判断。取证观测可以提高风险或请求复核，但不能独立证明篡改。语义模型不能推翻确定性的契约失败。

验证器记录服务方、模型版本、阈值、响应有效性与退化状态。必需服务方失效、强制证据不受支持或硬冲突会产出未知。回执承诺策略、证据、观测、决策、可容许转移与一个共享的不可变决策标识符。

\paragraph{可重构性作为两阶段原子事务。} 审计正确地指出：若“写回执”与“执行 $a^+$”是两个可分离的步骤，则存在一个时间窗——转移已发生而回执尚未落盘（或反之）——在此窗内 $\phi_R$ 的真值无从谈起。我们据此把授权—执行链规定为一个\emph{两阶段提交}：第一阶段\emph{预备}，守卫在冻结承诺上求值四条谓词并把完整回执写入持久日志，此时尚不执行 $a^+$；第二阶段\emph{提交}，当且仅当回执已持久化且 $\phi_R=\top$ 时，执行 $a^+$ 与“标记回执为已生效”作为一个\emph{原子}操作一并完成。任一阶段失败都回滚到“未授权”，绝不留下“已转移但不可重构”的中间态。因此 $\phi_R$ 不是一个事后检查，而是提交阶段的\emph{前置条件}：一个生产集成必须把验证、回执持久化与结算原子地绑定；参考验证器是一个一致性预言机，而非“每条部署路径都强制该守卫”的证明。我们承认这只把不可重构风险从应用逻辑推到了底层事务基元的正确性上，并在局限性中重访这一残余假设。

\subsection{一个 Worked Replay 案例}
\label{sec:worked}
为把守卫行为具体化——并展示它作用于一个强裁判会接受的输入——我们端到端地追踪一个 $D_B$（跨任务重放）样本包。该例仅作说明，取自基例模板族而非确证数据；不从中作任何结论。

\emph{设定。} 需求方发布任务 $q_{412}$：“在 18:00 UTC 前，从账户 \texttt{@acme} 发布句子 \texttt{Launch is live}。”契约创建时义务被冻结：标准版本 $v_3$、所需模态 \texttt{screenshot}、强制服务方 \texttt{web-archive}、决策方法 \texttt{selective}、失败行为 \texttt{abstain}。产出方随后提交一张截图 $e^\dagger$，它确实显示 \texttt{@acme} 上的 \texttt{Launch is live}——但该截图是为\emph{另一个}任务 $q_{197}$ 在标准版本 $v_2$ 下产生的，正被重放。所呈现的证据摘要是真实的；只有它的\emph{绑定}是错的。

\emph{裁判看到了什么。} 通用裁判 J0 与策略感知裁判 J1 读取证据后，都以高置信度报告 $\mathrm{success}$：句子、账户名与版式恰是请求所指名的。信息对齐裁判 J1-IM 额外以提示文本收到已承诺标识符（$q_{412}$、$v_3$）与捕获时记录的摘要（$q_{197}$、$v_2$），原则上\emph{可以}注意到两个任务标识符不同；它是否注意到，取决于模型是否去关注埋在上下文中的两个标识符，而非那张视觉上令人信服的截图。这恰是图~\ref{fig:indist} 的不可区分对：相同证据、相反真值、判别性事实只存在于上下文 $c$ 中。

\emph{契约做了什么。} 守卫在已承诺元数据上独立求值四条谓词：
\begin{itemize}
\item $\phi_P=\top$：生效标准是 $v_3$，与已承诺版本一致——无漂移。
\item $\phi_B=\bot$：绑定谓词以确定性相等比较把提交方的已承诺任务/版本 $(q_{197},v_2)$ 与契约的 $(q_{412},v_3)$ 相比；二者不同，故绑定\emph{失败}。
\item $\phi_V=\top$：选择性验证器仅读证据，确信句子存在。
\item $\phi_R=\top$：全部回执字段齐备。
\end{itemize}
由 $\bot$ 吸收，无论其余三个合取项如何满足，$\bigwedge_i\phi_i=\bot$，偏真守卫把 $\bot\mapsto$\emph{拒绝}（定理~\ref{thm:soundness}）。尽管语义内容正确、验证器也自信，正向动作仍被扣留：该转移之所以不安全，不是因为截图伪造，而是因为它\emph{未绑定到本任务}。回执记录了决定性合取项（$\phi_B=\bot$，连同两份摘要与两组标识符对），故审计者事后能精确重构“为何授权被拒”——这正是 J0/J1 无法提供的属性，因为它们从未表示过绑定。

这个单一案例同时展示了全文的每个活动部件：分离（相同证据、依上下文而定的真值）、“被提供信息的裁判”与“确定性比较”之间承重的区别、$\bot$ 吸收压过一个自信的 $\top$，以及可重构性把这次拒绝变成一份可审计记录。

\subsection{一个已部署的结算环路作为参考实例化}
\label{sec:instantiation}
前述原型是抽象描述的。为说明验证契约并非纸面构造，我们把六元组 $K=\langle I,P,B,V,A,R\rangle$ 映射到一个\emph{已独立部署、持续运行}的结算环路的真实组件上。该环路实现了一条完整的生产回路：\emph{活动发布 $\to$ 证据绑定 $\to$ 选择性验证 $\to$ 链上结算 $\to$ 回执}。为遵守双盲评审要求，本节刻意以\emph{去标识（de-identified）}方式描述该系统的能力，而\emph{不}指名平台、域名或合约地址；系统身份、致谢与链上制品指针留待评审通过后的定稿（camera-ready）揭示。\footnote{为维持双盲，本节所述已部署系统的名称、运营方与链上合约地址在评审阶段被隐去。一份包含系统标识、公开合约地址与只读审计入口的补充材料将在评审通过后提供；其内容哈希已随预注册一并加时间戳，以证明本节所述系统在投稿前即已存在并运行，而非事后补造。}

\begin{table}[t]
\centering
\footnotesize
\caption{把验证契约的六元组 $K=\langle I,P,B,V,A,R\rangle$ 逐字段映射到一个已部署结算环路的真实组件。系统以去标识方式描述；末列给出各字段被冻结或强制的时机，可与图~\ref{fig:guard} 的守卫及 \S\ref{sec:study} 的两阶段提交对照。}
\label{tab:instantiation}
\resizebox{\columnwidth}{!}{%
\begin{tabular}{@{}p{2.7cm}p{5.3cm}p{3.4cm}@{}}
\toprule
契约字段 & 已部署环路中的对应组件 & 冻结 / 强制时机 \\
\midrule
$I$ 声明结果与不可变约束 & 活动发布时锁定的结果声明与不可变参数（奖励、截止时间、允许的下游动作） & 活动发布即写入，此后不可改 \\
$P$ 事前证明义务 & 发布时冻结的验收标准版本与逐义务清单 & 首份证据被提交\emph{之前} \\
$B$ 证据绑定 & 把提交证据绑定到活动标识、标准版本、服务端回执时间与证据摘要 & 证据提交时刻 \\
$V$ 选择性验证器 & 多模态判断 $+$ 确定性规则/摘要检查，产出 $\top/\bot/?$；服务方失效或硬冲突产出 $?$ & 验证时刻 \\
$A$ 可容许转移 & 白名单化的链上结算 / 访问授予动作 & 提交阶段（两阶段提交） \\
$R$ 决策回执 & 每次决策产出的不可变回执与共享决策标识符（链上可查） & 与结算\emph{原子}绑定 \\
\bottomrule
\end{tabular}%
}
\end{table}

表~\ref{tab:instantiation} 表明该映射不是隐喻：该系统在活动发布时就冻结了 $I$ 与 $P$（对应事前一致性），在证据提交时刻计算 $B$ 的绑定摘要，在结算前经由选择性验证器求得 $V$，把 $A$ 限制在一组白名单动作上，并把 $R$ 与链上结算\emph{原子地}绑定——正是 \S\ref{sec:study} 所规定的两阶段提交。因此本文的抽象授权边界确有一个已在真实、不可逆结算路径上运行的对应物，而非仅存在于纸面。

\paragraph{为何这些机制只可能来自一个真实系统。} 上述去标识化描述之所以可信，是因为其中承重的性质并非任何纸面契约能事后补写，而只能来自一个必须直面\emph{不可逆结算}的运行系统所被迫做出的工程取舍。其一，绑定摘要 $B$ 由\emph{服务端}在收妥证据的瞬间计算，时间戳取自服务端权威时钟而非提交者自报，因此“证据属于本活动、本冻结标准”这一判断是防篡改的、可被第三方独立重算的，而非信任客户端的一面之词。其二，链上结算不可撤销，故系统在\emph{提交（commit）之前}即以一道守卫（图~\ref{fig:guard}）拦截：唯有 $V=\top$ 才进入结算，$V=?$（服务方失效或硬冲突）\emph{绝不}自动放款，而是挂起、托管或升级人工，$V=\bot$ 直接拒付——这正是 \S\ref{sec:study} 的两阶段提交在真实资金路径上的强制形态。其三，可容许转移集合 $A$ 被\emph{物理地}收敛为一组白名单动作，使任何一次错误裁决的“爆炸半径”由结算层本身界定，而非依赖裁判自律。其四，每笔决策产出的回执 $R$ 与其链上结算共享同一决策标识符并被\emph{原子地}绑定，因此可重构性不是一句承诺，而是一件业已存在、任何人皆可只读查询的链上制品。正是这些“先于不可逆动作强制、事后可独立重放”的约束，构成了一个纸面契约无法伪造、唯有已部署系统才会具备的证据。

\paragraph{作为确证数据的来源。} 组 B（外部结果）的确证数据拟\emph{从这个独立运行的生产结算系统采集}：真实活动产生真实的证据提交、结算，以及（在争议时）真实的申诉后果，其结果由\emph{对契约机制不知情的盲评者}标注。这与 \S\ref{sec:study} 及 \S\ref{sec:limitations} 的独立性论证一致：由于该系统独立于本文的标注流程运行、且盲评者不知晓协议内部，外部结果标签不会退化为协议自身的规则。它还部分弥合了 \S\ref{sec:limitations} 所坦承的缺口——在生产环路中，$V$ 所依赖的危害是\emph{被观测到的}真实结算/申诉后果，而非反事实重构。我们不掩饰潜在的利益冲突：由于该已部署系统与作者相关，我们正是通过对评定者隐藏系统身份、预测与结算状态（\S\ref{sec:limitations} 的盲评协议）来阻断“用自家系统证明自家系统好”的循环，并在定稿时完整披露该关系以供读者独立判断。我们重申 \S\ref{sec:limitations} 的告诫：对确定性缺陷类（$D_P,D_B$）锚点仍与规则共享输入，真实结算后果只对以 $V$ 为主的类提供真正独立的危害信号。在评审阶段，正文仅表述为“确证数据从一个独立运行的生产结算系统采集，由不知情盲评者标注”；该系统的身份与链上制品在定稿时补齐。

\section{JOVE-Core}
JOVE-Core 覆盖六类数字任务：社交内容、账户配置、内容发布、表单或文档提交、时间敏感的数字状态，以及跨证据一致性。物理执行不在确证范围之内。

独立单元是一个取得知情同意的基例，包含一个低风险请求、冻结契约、有效证据、溯源与放行边界。匹配变体在保留无关表面特征的同时，改动一个目标属性：事后标准漂移、跨任务或跨版本重放、服务方失效、完整性冲突，或不完整回执。同一基例的全部变体留在同一划分内，且没有任何基例同时进入校准分区与确证分区。

\paragraph{真值，及其独立性的诚实边界。} 为避免把成功定义为“与契约一致”——那会让一个执行契约的系统按构造取胜——确证标签是一个\emph{下游危害}，其\emph{实现}（是哪个系统、如何决策）对评定者隐藏。对每次转移，我们询问：放行该正向动作是否会（a）在不属于本任务或本冻结标准的证据上转移价值或授予访问，或（b）错误地拒绝一位合法提交者。评定者仅从请求、证据与生效标准判断，对系统身份、预测、变异类别、结算状态、协议内部细节以及彼此的评分均不知情。

我们坦率说明这种盲评达成与未达成之事。在研究者构造的、取得同意的低风险基例中，并无真实付款、真实访问授予或真实被拒的提交者；因此危害是一个关于\emph{未曾发生}后果的\emph{反事实}判断。对确定性缺陷类（$D_P$ 漂移、$D_B$ 重放），一位评定者判断“该证据是否属于本任务与本冻结标准”时，在信息意义上正是在手工重新执行事前谓词与绑定谓词。因此恰恰对这些类，危害标签对协议的\emph{实现}是盲的，却并非与协议的\emph{规则}在信息上正交：锚点与谓词共享其输入，没有任何标注协议能在没有真实下游后果可观测的情况下把二者分离。故我们\emph{不}宣称真值在一般意义上与规则无关。它的独立性只在概率性谓词（$V$）处才是真切的——那里危害取决于一件事实性事项（所宣称结果是否真的发生），评定者无需执行契约即可评估之——以及在真实结算或申诉后果被工具化的类上。对 $P$ 与 $B$，我们在明确告诫下报告结果：锚点是对同一规则的人工重算；我们把弥合这一缺口——通过引入真实结算与申诉后果，使危害被观测而非被重构——视为必要的未来工作，而非当前设计已达成之事。契约的“冻结义务下的充分性”仍是一个对照危害验证的\emph{预测器}，而非目标本身的定义；一个满足自身义务却仍产生危害的契约，将被计为对它不利。

\paragraph{样本量与裁定（多数决准入，一致性分层）。} 确证集以 200 个基例、至少 600 个匹配样本包为目标，其规模由第 7.5 节的功效分析确定，以在 80\% 功效下检出预注册的最小效应。每个样本包由 3 名评定者独立标注，危害标签按\emph{多数决}赋予。审计正确地指出：以“三名评定者全一致”作为准入判据会系统性地\emph{丢弃更难裁定的样本包}，而这些样本包很可能非随机地富集于某些缺陷类，从而在正式结果中引入选择偏差。我们据此\emph{放弃}把全一致准入作为主分析。主分析在\emph{全部}获得多数标签的样本包上进行——即准入判据为“存在多数”，$2$--$1$ 分歧的样本包按其多数标签\emph{纳入}而非留出。为了不丢失一致性信息，我们把每个样本包按其评定者一致性\emph{分层}（一致 vs.\ $2$--$1$），并在主分析之外预注册两项一致性敏感性分析：（a）仅在一致样本包上重估首要效应，（b）把 $2$--$1$ 样本包的标签按其不确定性加权。若主分析（多数决准入）与敏感性分析（一致准入）在效应方向上不一致，我们如实报告该分歧而不择优呈现。另外，仅作为语料级信度报告，我们在整个已标注集上计算 Krippendorff's $\alpha$ \cite{hayes2007answering}，并预注册 $\alpha\geq 0.67$ 作为整个语料的接受门槛；$\alpha$ 从不用于准入或排除单个样本包。由于评定者一致性本身可能与缺陷类相关（第~\ref{sec:study} 节），我们\emph{按缺陷类}报告准入率、一致性率与确证效应，使更难裁定的类不会被悄然从正式结果所关涉的类中移除。不可约的标准歧义被记录为一个带标签的类别并报告其占比，而非被悄然排除，以免排除抬高效应。

\section{实验设计}
\label{sec:study}
\subsection{系统}
J0 是一个通用的事后裁判。J1 是一个强策略感知裁判。J1$+$abstain 用一个置信度阈值化的弃权接口增强 J1，使覆盖率可被调节并与契约匹配。J1-IM 是信息对齐裁判：它收到与 K-Full 验证器\emph{完全相同的输入字段}，但不含对这些字段的确定性比较裁决（见下文逐字段构造）。J1-DT 是\emph{确定性工具增强}裁判：在 J1-IM 的全部字段之上，它\emph{额外}被授予调用一组确定性工具的能力——一个标识符/版本相等检查器与一个摘要相等检查器——并被明确提示在裁决前可以调用它们。K-P、K-B、K-S 增量地加入事前策略、绑定与选择性。K-Full 实现全部四条属性。每个系统都报告调用次数、token、时延与货币成本。

\paragraph{为何需要 J1-DT，以及它检验什么。} 单有 J1-IM 会招致一条公正性反驳：它必须靠概率模型对上下文中两串字符的 token 级注意力去\emph{注意到}不匹配，而 K-Full 靠一个确定性相等算子去\emph{强制}它；若 K-Full 胜出，怀疑者可归因于“K-Full 有工具而基线没有”，而非归因于结构性强制本身。J1-DT 关闭这条退路：它\emph{拥有}与 K-Full 相同的确定性检查能力，差异仅在于这些检查是\emph{可选的、由模型自行决定何时调用}，而非被守卫\emph{无条件强制并以 $\bot$ 吸收}。于是 K-Full 对阵 J1-IM 度量“提供确定性比较的价值”，而 K-Full 对阵 J1-DT 度量“把该比较设为强制结构、而非交由模型裁量”的\emph{增量}价值——这才是本文真正主张的那一部分。若 J1-DT 可靠地自行调用工具并达到 K-Full 的结构一致性，则本文对“强制优于裁量”的主张被削弱，我们据实报告（第~\ref{sec:results} 节，条件 (2$'$)）。

\paragraph{J1-IM 的逐字段构造，以及“信息对齐”的确切含义。} 因为“信息 vs.\ 结构”正是本文所争的界线，J1-IM 的输入被逐字段规定并在数据之前冻结。J1-IM 在其提示中收到 K-Full 验证器所消费的每一个\emph{输入}——原始证据、冻结标准文本、声明的义务、已承诺的任务与版本标识符、已承诺与所呈现的证据摘要、以及服务方退化状态——皆以文本呈现。它\emph{不}收到任何比较的\emph{结果}：所呈现摘要是否等于已承诺摘要、版本是否匹配、弃权决策，或回执。这些结果正是被测试的结构性强制。

我们在此精确说明什么被匹配、什么未被匹配，因为这是整个设计的关键。J1-IM 在决策可得的原始信息上被匹配：它被交予完全相同的字段，包括两份摘要与两个版本标识符，原则上能报告不匹配。它\emph{不}被交予对这些字段的确定性相等裁决——那正是 K-Full 的贡献。这是“信息对齐”唯一公平的读法：输入相同，差异仅在于不匹配是被一个可靠的确定性比较捕获、还是必须由一个读取上下文中两串字符的概率模型去注意。我们不宣称一个文本 LLM 会重算哈希；要点恰相反——绑定之所以有价值，正因为相等检查被委派给一个确定性算子、而非留给模型的 token 级读取，确证对比度量的正是这一委派的价值。把比较裁决本身交给 J1-IM 会使它变成 K-Full，故我们不这样做；不给它输入字段会使它变成仅依赖证据、并由定理~\ref{thm:separation} 平凡落败，故我们也不这样做。该基线被刻意置于二者之间，其提示中的每个字段都被记录在预注册字段映射中，使它不能被悄然增强或削弱。

\subsection{首要端点：结构一致性与外部结果的分离}
审计指出，把四条属性的效果压进单一“不安全正向转移率”会把两类在认识论上不同的宣称混在一起：其一，守卫是否\emph{结构性地}执行了它所承诺的确定性义务（$P$ 事前一致、$B$ 绑定、$R$ 可重构）；其二，守卫是否真的降低了一个\emph{外部的、独立评定的}下游危害（其判别性事实是概率性的，主要经由 $V$）。前者的真值可由确定性重算无歧义地确定，后者则依赖盲评危害标签。我们据此把确证端点\emph{拆成两组}，并分别预注册。

\paragraph{端点组 A（结构一致性，确定性）。} 对每条确定性谓词 $i\in\{P,B,R\}$，端点是\emph{结构违反漏放率}——即一个带匹配缺陷 $D_i$ 的样本包被错误放行 $a^+$ 的比例。其真值由对冻结承诺的确定性重算给出，不经盲评，因而是一个关于\emph{机制是否按定义运作}的对齐检查，而非外部有效性宣称。我们预期 K-Full 在这些端点上按构造达到 0（定理~\ref{thm:necessity} 的经验镜像），并把它作为\emph{可证伪的实现正确性检查}呈现：任何非零值都暴露实现缺陷，而非验证抽象价值。

\paragraph{端点组 B（外部结果，独立评定）。} 首要\emph{科学}端点是不安全正向转移率——一个独立评定的下游危害，定义于第 7 节——作为在跨系统匹配的覆盖率水平上的\emph{成对比较}来评估。此端点的判别性事实主要落在概率性谓词 $V$ 所处理的情形，其真值由对协议实现盲的评定者给出，故承载本文的外部有效性宣称。首要对比是 K-Full 对阵 J1-IM，因为它固定输入信息、只改变强制是否为结构性：这是最直接隔离本文论点的对比。但它并非唯一的证伪路径；预注册的全部证伪条件枚举于第~\ref{sec:results} 节，其中任一成立即反驳核心宣称。K-Full 对阵 J1$+$abstain 是一个次要对比，仅在风险—覆盖率曲线有良好定义处报告（见下文告诫）；此前固定的“$\geq 70\%$”点仅作为敏感性分析的一个带标签切片保留，因为 J1 没有原生覆盖率控制、强加一个本身即是一种设计选择。

我们坚持组 A 与组 B 的分离，因为把二者相加会让一个\emph{结构上}正确却在\emph{外部结果}上无益的守卫，靠 $P/B/R$ 上的确定性零来虚高其总分。仅组 B 可用于宣称验证契约相对信息对齐基线的经验优越性。

\paragraph{关于基于曲线之比较的告诫。} K-Full 与 J1$+$abstain 不共享一个公共覆盖率机制：J1$+$abstain 随其置信度阈值扫描而画出一条平滑曲线，而 K-Full 的确定性谓词（$P,B,R$）是不随覆盖率旋钮变化的硬门，故其选择性行为只经由概率性谓词 $V$ 变化。因此我们把二者视为\emph{不}处于公共覆盖率轴上，且\emph{不}把它们之间单一的成对 AURC 差当作有良好定义来报告。取而代之，我们只在明确匹配的接受覆盖率点上、且限于 $V$（因而覆盖率扫描）为起作用谓词的缺陷类上，比较 K-Full 与 J1$+$abstain；AURC 逐系统作为系统内摘要报告，而非作为跨系统成对面积差。

\paragraph{危害不是一个标量：$H_{\mathrm{FA}}$、$H_{\mathrm{FR}}$、$H_{\mathrm{AUDIT}}$。} 审计正确地指出，把“危害”当作单一标量会把方向相反的错误混为一谈。我们据此把结果危害显式分解为三种\emph{独立}成分，各自单列报告、绝不先验相加：（a）\emph{误接受危害} $H_{\mathrm{FA}}$——放行了一个本应被拦截的不安全正向转移（即不安全正向转移率，本文的首要科学端点）；（b）\emph{误拒绝危害} $H_{\mathrm{FR}}$——拦截或弃权了一个本应放行的安全转移，其代价是可用性与吞吐损失；（c）\emph{审计危害} $H_{\mathrm{AUDIT}}$——即便转移决策正确，回执仍无法支撑事后重构（缺失决定性义务、证据承诺或退化路径）。三者在不同的损失尺度上度量、服务不同的利益相关方，把它们相加会掩盖一个“靠制造大量 $H_{\mathrm{FR}}$ 来压低 $H_{\mathrm{FA}}$”的退化解。仅当明确声明一个损失矩阵时我们才做加权组合，且该矩阵作为敏感性维度扫描（见下）。

\paragraph{校准的可操作化。} 我们不空谈“校准”，而预注册其度量与阈值：$C$ 层输出的分数以期望校准误差（ECE，15 个等频箱）\cite{guo2017calibration} 与可靠性图独立评估，并报告 Brier 分数分解为可靠性/分辨率/不确定性三项；当采用保形集时，报告其在预注册 $\alpha$ 处的经验覆盖率与平均集合大小 \cite{vovk2005algorithmic,angelopoulos2023conformal}。校准诊断在\emph{校准划分}上估计、在\emph{确证划分}上核验，二者由防泄漏门隔离；校准质量作为一个\emph{独立}端点报告，绝不折叠进不安全转移率。

我们逐系统报告风险—覆盖率曲线下面积（AURC）作为系统内摘要，并在跨系统匹配的覆盖率水平上报告成对风险差、风险比与 95\% 基例聚类自助区间（10{,}000 次重采样）\cite{efron1994bootstrap}。为与第 7.5 节的分层功效核算及按缺陷类的准入报告保持一致，自助\emph{按缺陷类分层}并\emph{在基例层聚类}重采样，使区间同时反映类构成与基例内相关，而非把所有样本包当作独立同分布。关键地，头条效应被报告为一个\emph{逐缺陷类的条件率} $P(\text{危害}\mid\text{缺陷类})$，因为聚合不安全率是作者所选缺陷占比的函数、不是一个可解释的效应量。我们额外报告总自动化覆盖率、接受覆盖率与每负例的不安全接受，使增益不能靠压低覆盖率来制造。确证决策以一个\emph{对称}的 1:1 误接受:误拒绝损失作为首要工作点，使任何系统都不因一个恰好奖励其偏好错误方向的成本不对称而占优；一个只靠更多拒绝或弃权来降低不安全转移的系统，在等成本下仍须显示净增益。2:1、5:1、10:1 的非对称损失——其下一次错误的价值转移或访问授予被视为逐步比一次可恢复的错误拒绝更昂贵——作为敏感性分析报告，我们并报告系统排名在各比率间是稳定还是在某处反转。

\subsection{机制与泛化}
属性消融度量标准漂移转移、重放接受、声明弃权条件下的强制转移，以及审计重构。至少两个任务类与一个变异构造源被留出。一个在 J1-IM 下消失、或只在已见模板或已见变异源上成立的增益，不被归因于契约抽象，而被归因于信息或记忆。

\subsection{人类研究}
复核研究用三臂而非两臂，以把信息组织与契约结构分离：（i）原始任务与证据；（ii）一个信息对齐包，以有组织但非契约的格式呈现相同事实；（iii）结构化契约包，三者都略去最终的机器裁决。要把任何效应归因于契约抽象，需要第（iii）臂胜过第（ii）臂，而不仅是胜过第（i）臂。首要人类结果是不安全接受；次要结果包括准确率、时间、置信度、工作负荷与不当依赖。一个单独的审计研究比较规范回执与普通日志，度量对策略、决定性义务、证据承诺、退化路径与转移资格的重构。

\subsection{样本量与功效}
确证集以 200 个基例、至少 600 个匹配样本包为目标，每个样本包 3 名评定者。以下功效数字是一个\emph{规划假设}，而非一项已达成的保证：其输入（基线率、最小效应、组内相关）均为\emph{文献先验 / 试点前的假定值}，尚未由试点数据估计，故 $n=600$ 应读作一个\emph{设计目标}而非“已获 80\% 功效”的宣称。在此规划假设下——某缺陷类内 J1-IM 的基线不安全率 0.30、预注册最小可检出效应为绝对值 0.10 的下降（即降至 0.20）、基例聚类的组内相关 0.2——冻结的成对聚类自助模拟（\texttt{scripts/power-paired-cluster.mjs}，随机种子 20260809、3000 次重复、每基例 3 个变体、双侧 $\alpha=0.05$）给出：\emph{在单一缺陷类内}需 110 个基例（即 330 个样本包）方达 $0.808$ 的规划功效。这一点估计澄清了此前"约 80\%"的模糊表述：功效是\emph{在类内}定义的，而非在合并的 600 样本包上；600 这一目标之所以显著高于单类所需的 330，恰恰来自下述三项规划余量，而非冗余。我们据审计意见补入三项此前缺失的规划要素：（i）\emph{分层}——功效按缺陷类分层核算，因为效应与准入率均可能随类变化，我们对预注册的\emph{最关键缺陷类}保证其类内基例数达到上述 110 的规模并单独报告其类内可检出效应，其余类的类内效应作为探索性结果报告而不参与受保护的确证宣称；（ii）\emph{稳健性}——冻结的敏感性扫描（\texttt{scripts/power-sensitivity.mjs}）在基线 $\in\{0.25,0.30,0.35\}$、效应 $\geq 0.10$、组内相关 $\in\{0.1,0.2,0.35\}$ 的整个合理参数域内，单类所需基例数介于 70 与 140 之间，故预注册的 200 基例总量对该域内的最坏情形（$0.35\!\to\!0.25$、$\rho=0.35$ 需 140 基例）亦保有余量；（iii）\emph{损耗}——我们假定同意撤回、评定者流失与标准歧义会带来约 15\% 的样本损耗，据此对招募目标\emph{上浮}，使净样本仍达 600；（iv）\emph{多重性}——由于我们预注册了多个端点（结构组 A 的 $P/B/R$、外部组 B 的首要与次要对比、以及 J1-DT 对比 (2$'$)），我们以 Benjamini–Hochberg 控制外部结果族的错误发现率，并把首要对比（K-Full vs.\ J1-IM）作为受保护的主检验先行。上述基线率与组内相关为\emph{规划先验}，将由试点按预注册的\emph{先验替换协议}以其\emph{对功效更不利}的 95\% 自助边界替换后重跑脚本冻结；而最小可检出效应 0.10 是由"多大的安全改善才值得宣称"这一\emph{决策}锚定的量，\emph{不}由任何数据（试点或确证）重估，以免对确证检验偷看。试点估计只能使冻结样本量\emph{持平或上调}，指向更小样本量时仍以 200 基例／600 样本包为下限。若招募得到的同意案例更少，或试点估计显著偏离上述先验，研究报告\emph{实际}达成的功效并把确证宣称降级为探索性，而非放松阈值。注册条目、冻结的分析脚本及其提交哈希在任何确证数据被采集之前即加时间戳。

\section{结果与证伪}
\label{sec:results}
本节必须在取得同意的确证数据被独立标注后、由冻结脚本生成。目前没有可报告的确证结果。

为在数据采集之前验证分析接线、表格布局与指标管路，我们在一个\emph{合成排练}语料上运行冻结脚本，并核验每个证伪条件都能被端到端计算。\textbf{但我们刻意不把排练所得的任何合成数值印入正文}：表~\ref{tab:mainresults} 与~\ref{tab:propertyresults} 仅展示\emph{结构模板}，全部单元格为 \texttt{--} 占位符，以杜绝合成数字（即便标注为“说明性”）被下游读者、二次引用或语言模型误当作经验结果。排练数字仅存活于被合成隔离门（自第~\ref{sec:results} 节起，并在“可复现性”中经\emph{路径级}隔离强制）隔离的夹具文件（\texttt{results/mock-metrics.json}）中，绝不进入确证路径，并在提交前被冻结的确证指标覆盖。

% JOVE_RESULTS_BLOCK_BEGIN
\begin{table*}[t]
\centering
\caption{主结果的\emph{结构模板}。所有单元格显示为 \texttt{--} 占位符：本协议快照下\emph{没有}可报告的确证数字，我们刻意不填入任何合成数值，以杜绝被误读为经验结果。真实值将在取得同意的数据被独立标注后由冻结脚本填入。新增的 J1-DT 行为确定性工具增强裁判基线。UPT 表示不安全正向转移率（即 $H_{\mathrm{FA}}$）；Coverage 为自动化覆盖率；AURC 与 ECE 为系统内摘要。}
\label{tab:mainresults}
\resizebox{\textwidth}{!}{%
\begin{tabular}{lrrrrrr}
\toprule
系统 & UPT $\downarrow$ & 覆盖率 $\uparrow$ & AURC $\downarrow$ & ECE $\downarrow$ & 成本 & 时延 \\
\midrule
J0 通用裁判        & -- & -- & -- & -- & -- & -- \\
J1 策略感知        & -- & -- & -- & -- & -- & -- \\
J1-IM 信息对齐     & -- & -- & -- & -- & -- & -- \\
J1-DT 确定性工具   & -- & -- & -- & -- & -- & -- \\
K-P 策略契约       & -- & -- & -- & -- & -- & -- \\
K-B 绑定二值       & -- & -- & -- & -- & -- & -- \\
K-S 选择性         & -- & -- & -- & -- & -- & -- \\
K-Full             & -- & -- & -- & -- & -- & -- \\
\bottomrule
\end{tabular}%
}
\end{table*}
% JOVE_RESULTS_BLOCK_END

\begin{table}[t]
\centering
\caption{属性级结构模板。所有单元格显示为 \texttt{--} 占位符；本协议快照下无任何经验数字，真实值将由冻结脚本在确证数据到达后填入。}
\label{tab:propertyresults}
\resizebox{\columnwidth}{!}{%
\begin{tabular}{lrrr}
\toprule
属性测试 & 基线 & 契约变体 & 成对差 \\
\midrule
标准漂移转移 & -- & -- & -- \\
重放接受     & -- & -- & -- \\
强制转移     & -- & -- & -- \\
审计重构     & -- & -- & -- \\
\bottomrule
\end{tabular}%
}
\end{table}

核心宣称在下列任一预注册条件下失败；每条各自独立，任一成立即反驳该宣称。(1) 在匹配覆盖率下，K-Full 未能相对 J1-IM 降低逐缺陷类的不安全率，且成对 95\% 区间跨越零。(2) K-Full 相对 J1-IM 的优势消失——即效应由信息而非结构解释。(2$'$) K-Full 相对 J1-DT（确定性工具增强裁判）的优势消失——即效应由“持有确定性检查”而非“把检查设为强制结构”解释；此时“强制优于裁量”的主张被反驳。(3) 在覆盖率扫描有良好定义的缺陷类（谓词 $V$）、匹配接受覆盖率下，K-Full 未能改进 J1$+$abstain。(4) 增益不能迁移到留出任务类或留出变异源；由于确定性协议无法记忆模板，条件 (4) 检验的是\emph{习得基线}的泛化与效应归因，而非 K-Full 的记忆。(5) 不安全转移的下降伴随接受覆盖率的下降，大到每负例不安全接受并未改善——操作化为在任何匹配覆盖率水平上均无改善，替代此前未定义的“主要”。(6) 规范回执未能比普通日志改进独立重构。若仅 K-P 即解释全部效应，则后续属性不被宣称为必要。(7) 属性消融是对模型内必要性结果（定理~\ref{thm:necessity}）的经验探针，我们诚实陈述其范围：对确定性谓词（$P,B,R$），定理的见证按构造恢复匹配危害，故其消融确认的是接线而非提供一个独立检验，我们不把它当作后者呈现。承载经验分量的消融是概率性谓词 $V$——移除它是否在真实数据上重开匹配危害，是真正不确定的；若该危害未出现，则 $V$ 的必要性不外推到缺陷模型之外，我们据实报告而非宣称经验最小性。若人类研究的第 (iii) 臂未胜过第 (ii) 臂，则人类收益被归因于组织而非契约。若结构化人类复核加剧了自动化偏见，该负面结果仍保留在文中。

\section{预答反驳}
\label{sec:objections}
我们陈述一位批判性读者最先提出的四条反驳，并针对本文自身的宣称、而非更弱的转述，逐一作答。

\paragraph{“这不就是加了一堆术语的元数据校验吗？”} 若四条谓词只是字段检查，本文便是一次命名练习。真正承重的内容不在于我们检查任务标识符，而在于\emph{做这件事的正确性被锚定在何处}。分离（定理~\ref{thm:separation}）表明一个容量无界的仅依赖证据的裁判无法闭合一个不可区分对；必要性（定理~\ref{thm:necessity}）表明每删去一条谓词都会重新引入一个特定的匹配危害；评估随后把真值锚定到\emph{外部下游危害}、而非“字段是否匹配”（\S7）。一个纯元数据校验器一样都没有：它把成功定义为自身规则，这正是我们盲评危害标签所要避免的同义反复。被测宣称之所以可证伪，恰因为一个满足自身义务却仍产生危害的元数据校验器会被计为对我们不利（\S\ref{sec:results}，条件 5）。

\paragraph{“为何不用 TEE 让绑定不可伪造？”} 一个飞地能保证裁判\emph{未经修改地运行}、输入\emph{被完整接收}；它不能判定一个完整、真实的证据在一条上周就冻结的标准下\emph{是否有资格}授权本次转移。\S3 的全部四类缺陷都能挺过完美的飞地证明：一张被重放的截图是真实捕获的，一个被证明的裁判仍读到一条漂移了的标准，一个强制服务方仍是宕机的。因此 TEE 加固的是喂给 $V$ 与 $B$ 的\emph{证据通道}；它并不提供 $P$、选择性弃权，或一份危害锚定的回执。我们把飞地视为一个加固输入——并注意它们自身在瞬态执行攻击下的机密性失效 \cite{vanbulck2018foreshadow}，这是又一个不把证明当作\emph{那条}授权边界的理由。

\paragraph{“一个强 LLM 不就能读两个哈希与两个任务 ID、抓出不匹配吗？”} 有时能——而这正是信息对齐基线 J1-IM \emph{逐字收到两份摘要与两个标识符}的原因（\S\ref{sec:study}）。我们不宣称语言模型无力注意到不匹配；我们宣称的是：把绑定检查留给模型对上下文的 token 级注意力，相对一个确定性相等算子而言是不可靠的，而确证的 K-Full 对阵 J1-IM 对比度量的正是这一缺口在真实数据上的大小。若缺口为零——若模型可靠地抓住了确定性检查所抓住的东西——则 \S\ref{sec:results} 的条件 (2) 触发、核心宣称被反驳。于是本文的论点恰恰押注在这条反驳所提议的那个实验上。

\paragraph{“没有确证数据，这为何是一篇论文而非一份提案？”} 三项贡献中有两项不依赖待采数据。分离与必要性结果是定理，其真假独立于任何语料；评估仪器——预注册端点、对称损失工作点、逐缺陷类条件率、七个独立证伪条件，以及使“信息 vs.\ 结构”可检验的逐字段 J1-IM 构造——本身即是一项贡献，其充分性现在即可评判。数据将裁定的是经验量级，我们刻意拒绝预览它：合成排练（\S\ref{sec:results}）的存在只为证明管线能计算每个证伪条件，并由合成隔离门与确证路径隔开。

\section{局限与伦理}
\label{sec:limitations}
一个验证契约可以稳固一条不正确、歧视性或不正当的策略。协议资格不等于法律、社会或道德上的正当性。哈希绑定不证明捕获真实性。取证是不完整的。模型可能共享相关误差。人类复核者可能有偏或合谋。

我们的威胁模型容许一个自适应、对抗性的产出方，但确证数据由研究者构造的匹配缺陷组成。因此我们只检验对所预期并实例化的攻击类的鲁棒性，而非对一个针对已部署契约进行优化的自适应对手。相应地，确证宣称被限定于六类取得同意的数字任务、低风险合作情形，以及所枚举的缺陷族；泛化只对两个留出类与一个留出变异源评估，我们不宣称对未见攻击策略、或对物理的、高价值的、非合作的结果的泛化。

本研究排除法律、医疗、雇佣、信贷、生物特征、私密及隐蔽监控类决策。参与需要明确同意、有界留存、最小证据负担与一条申诉路径。原始敏感证据保持访问受控；公开证据默认以承诺、版本与派生观测呈现。

\section{可复现性}
该制品沿一条单向流水线组织，把每个阶段的产物冻结为下一阶段的只读输入，并对每个阶段的产物记录内容哈希，使对上游的任何回溯改动都\emph{可被检测}（而非在物理上不可能）：同意登记 $\to$ 证据捕获 $\to$ 盲评包构造 $\to$ 独立标注 $\to$ 系统运行与输出 $\to$ 评分 $\to$ 基例聚类自助分析 $\to$ 论文与图表生成。每个阶段以带模式（schema）校验的 JSON 记录落盘，原始评分以只追加方式保存，更正只能以附加、并附裁定文档的方式进行，绝不静默覆盖；因此任何事后改动都会破坏已记录的哈希链而暴露——正是该基准自身所研究的可审计性的镜像。我们不宣称加密意义上的防篡改：一个握有全部写权限的对手仍可重算整条哈希链，我们保证的是\emph{篡改可检测}，而非篡改不可能。

自动化门在提交前对制品逐项强制以下不变式，任一失败即阻断论文组装：（i）\emph{划分泄漏}检查，拒绝任何基例的变体跨越校准与确证分区；（ii）\emph{变异标签泄漏}检查，确保变异类别对评定者与系统均不可见；（iii）\emph{隐私审计}，度量被请求字段、被提交字段、含 PII 字段、留存字节与决策无关字段，并强制公开层只含承诺、版本与派生观测；（iv）\emph{匿名性}检查，确保盲评包不泄露系统身份或评定者彼此的评分；（v）\emph{回执语义}检查，验证每份回执携带策略、决定性义务、证据承诺、退化路径与转移资格；（vi）\emph{合成隔离}门，通过\emph{路径级}隔离——排练数据写入独立命名空间、合成指标显式标记、确证路径下不得出现任何合成或模板来源的文件——来阻止带合成或模板来源的行进入确证资产。该门是文件与命名空间层面的机械检查，并不对内容作语义判定。

提交要求一份取得同意的确证清单、封存的独立标签、不可变的系统输出、由冻结脚本生成的真实指标与自助制品、溯源哈希、重新生成的表与图，以及一次干净环境下的复现。复现包由 \texttt{build-reproduction-bundle.mjs} 产出，其 \texttt{MANIFEST} 的 sha256 稳定可校验，并在 \texttt{bundle\_type} 字段标注其为\emph{仅方法}包（确证数据尚未采集时）抑或\emph{含结果}包；\texttt{provenance.json} 保证确证行不携带任何模板或合成来源。发表时我们\emph{计划}把冻结的制品归档至一个带持久标识符（DOI）的公共仓库，使被弃用的发布仍可检索、既往结果保持可复现；本协议快照下该归档尚未创建。

一次合成排练可以验证代码路径、表格布局与指标管路能端到端运行，但\emph{不能}满足任何经验门：合成隔离门在任何确证数据被采集之前即在文件路径与命名空间层面把排练数字与确证路径隔开，并在提交时由冻结的确证指标覆盖。

\section{结论}
开放世界智能体的可靠性不能止于一个模型\emph{相信}某任务成功了。当概率性观测触发确定性后果时，一个系统必须陈述它为何有资格行动。验证契约把这一资格分解为事前规则、任务绑定的证据、选择性决策与可重构转移。这一抽象是否降低不安全转移，是一个由独立、可证伪的评估来回答的经验问题——而非由架构复杂度或工程完整度来回答。

\bibliographystyle{plainnat}
\bibliography{references}
\end{document}
